\begin{table}[htbp]
    \centering
    \small
    \setlength{\tabcolsep}{3.5pt}
    \resizebox{\textwidth}{!}{%
    \begin{tabular}{llcccccc}
        \toprule
        Architecture & Method & Del. advantage $\uparrow$ & Ins. advantage $\uparrow$ & Mean Sens.-$n$ (PCC) $\uparrow$ & Infidelity $\downarrow$ & Completeness $\downarrow$ & Seconds/sample $\downarrow$ \\
        \midrule
        Standard & Saliency & 0.356 $\pm$ 0.045 & 0.232 $\pm$ 0.028 & 0.028 $\pm$ 0.036 & 1.000 $\pm$ 0.000 & 1.00e+00 & \textbf{0.000305} \\
        Standard & Integrated Gradients & \textbf{0.396 $\pm$ 0.036} & \textbf{0.246 $\pm$ 0.029} & \textbf{0.739 $\pm$ 0.044} & \textbf{0.180 $\pm$ 0.074} & \textbf{7.80e-03} & 0.006874 \\
        Standard & SmoothGrad & 0.356 $\pm$ 0.044 & 0.233 $\pm$ 0.032 & 0.026 $\pm$ 0.038 & 1.000 $\pm$ 0.000 & 1.00e+00 & 0.003153 \\
        \midrule
        B-cos & Intrinsic B-cos & 0.221 $\pm$ 0.028 & \textbf{0.187 $\pm$ 0.009} & \textbf{0.422 $\pm$ 0.067} & \textbf{0.453 $\pm$ 0.070} & \textbf{1.62e-07} & \textbf{0.000434} \\
        B-cos & Saliency & 0.188 $\pm$ 0.032 & 0.174 $\pm$ 0.011 & 0.054 $\pm$ 0.013 & 1.000 $\pm$ 0.000 & 1.00e+00 & 0.000613 \\
        B-cos & Integrated Gradients & \textbf{0.223 $\pm$ 0.022} & 0.172 $\pm$ 0.016 & 0.350 $\pm$ 0.057 & 1.733 $\pm$ 0.108 & 6.58e-01 & 0.014209 \\
        B-cos & SmoothGrad & 0.187 $\pm$ 0.032 & 0.174 $\pm$ 0.013 & 0.048 $\pm$ 0.007 & 1.000 $\pm$ 0.000 & 1.00e+00 & 0.004439 \\
        \bottomrule
    \end{tabular}%
    }
    \caption{RQ2 explanation results. Deletion and insertion are AUC advantages over equal-area random perturbations. The adapted Sensitivity-$n$ summary is the mean of Pearson correlation coefficients (PCCs) computed separately at subset fractions from 10\% to 90\%. Values for deletion, insertion, mean Sensitivity-$n$, and infidelity are means $\pm$ half-widths of the 95\% cluster-bootstrap intervals. Higher is better for deletion, insertion, and mean Sensitivity-$n$; lower is better for infidelity, completeness error, and runtime.}
    \label{tab:rq2_explanation_methods}
\end{table}
